\section{Results}
\label{sec:results}

\subsection{Reference Prices}
\label{sec:ref_prices}

Table~\ref{tab:ref_prices} presents estimates of the effect of urgent procurement on log reference prices---the maximum price the government is willing to pay, set during the internal planning phase.

\begin{table}[ht]
  \centering
  \caption{Reference Prices}
  \label{tab:ref_prices}
  \small
  \begin{threeparttable}
  \begin{tabular}{lcccc}
    \hline\hline
    & (1) & (2) & (3) & (4) \\
    \hline
    Urgent Purchase & 0.164*** & 0.156*** & 0.027* & 0.027* \\
     & (0.056) & (0.050) & (0.014) & (0.014) \\[3pt]
    \hline
    Item FE & Yes & Yes & Yes & Yes \\
    Year FE & No & Yes & Yes & No \\
    Year-Month FE & No & No & No & Yes \\
    PBU FE & No & No & Yes & Yes \\
    Observations & 196,899 & 196,899 & 196,896 & 196,896 \\
    Within R$^2$ & 0.003 & 0.003 & 0.000 & 0.000 \\
    \hline\hline
  \end{tabular}
  \begin{tablenotes}
    \small
    \item \textit{Notes:} Dependent variable: log reference price. Sample: winners only, items with both ordinary and litigated purchases. Winsorized at 1\%/99\%. Standard errors clustered at the PBU level in parentheses. *** \textit{p}$<$0.01, ** \textit{p}$<$0.05, * \textit{p}$<$0.1.
  \end{tablenotes}
  \end{threeparttable}
\end{table}


The coefficient on Urgent Purchase is positive and statistically significant across all specifications. In the most parsimonious specification with item FE only (column~1), urgent purchases carry reference prices that are 17.8\% higher ($e^{0.164}-1$). Adding year FE (column~2) yields a similar estimate of 16.9\%. The magnitude decreases substantially in our preferred specification with item + year + PBU FE (column~3), where the premium is 2.7\% ($e^{0.027}-1$, $p < 0.10$), and remains stable at 2.7\% with year-month FE (column~4).

The attenuation from columns~(1)--(2) to columns~(3)--(4) is informative. It indicates that a substantial share of the raw price premium in urgent procurement is absorbed by PBU fixed effects, suggesting that certain buying units systematically handle more urgent purchases and also face higher prices for unrelated reasons. Our preferred estimate of 2.7\% represents the within-item, within-buyer, within-year premium attributable to urgency.

\subsection{Quantities}
\label{sec:quantities}

Table~\ref{tab:quantities} reports estimates for log quantities demanded.

\begin{table}[ht]
  \centering
  \caption{Quantities}
  \label{tab:quantities}
  \small
  \begin{threeparttable}
  \begin{tabular}{lcccc}
    \hline\hline
    & (1) & (2) & (3) & (4) \\
    \hline
    Urgent Purchase & -0.264 & -0.279 & -0.058 & -0.065 \\
     & (0.183) & (0.181) & (0.056) & (0.058) \\[3pt]
    \hline
    Item FE & Yes & Yes & Yes & Yes \\
    Year FE & No & Yes & Yes & No \\
    Year-Month FE & No & No & No & Yes \\
    PBU FE & No & No & Yes & Yes \\
    Observations & 196,899 & 196,899 & 196,896 & 196,896 \\
    Within R$^2$ & 0.003 & 0.003 & 0.000 & 0.000 \\
    \hline\hline
  \end{tabular}
  \begin{tablenotes}
    \small
    \item \textit{Notes:} Dependent variable: log quantity. Sample: winners only. Winsorized at 1\%/99\%. Standard errors clustered at the PBU level in parentheses. *** \textit{p}$<$0.01, ** \textit{p}$<$0.05, * \textit{p}$<$0.1.
  \end{tablenotes}
  \end{threeparttable}
\end{table}


The point estimates are negative across all specifications, consistent with reduced quantities under urgency, but are not statistically significant at conventional levels in our preferred specifications (columns~3--4). The magnitude of $-0.058$ in column~(3) corresponds to approximately 5.6\% fewer units. The imprecision likely reflects the high variance in quantities across purchase orders and the fact that quantities are partially determined by the nature of the court order itself (e.g., individual prescriptions specify specific amounts).

\subsection{Negotiated Prices}
\label{sec:neg_prices}

Table~\ref{tab:neg_prices} presents the core results on negotiated prices. Panel~A reports the total effect; Panel~B controls for log quantity to isolate the direct effect net of the quantity channel.

\begin{table}[ht]
  \centering
  \caption{Negotiated Prices}
  \label{tab:neg_prices}
  \small
  \begin{threeparttable}
  \begin{tabular}{lcccc}
    \hline\hline
    & (1) & (2) & (3) & (4) \\
    \hline
    \multicolumn{5}{l}{\textit{Panel~A: Total Effect}} \\[3pt]
    Urgent Purchase & 0.159*** & 0.151*** & 0.053*** & 0.056*** \\
     & (0.042) & (0.038) & (0.016) & (0.016) \\[3pt]
    \hline
    \multicolumn{5}{l}{\textit{Panel~B: Direct Effect}} \\[3pt]
    Urgent Purchase & 0.075 & 0.060 & 0.030* & 0.030* \\
     & (0.045) & (0.042) & (0.017) & (0.017) \\[3pt]
    Log Quantity & -0.321*** & -0.325*** & -0.392*** & -0.393*** \\
     & (0.026) & (0.026) & (0.029) & (0.029) \\[3pt]
    \hline
    Item FE & Yes & Yes & Yes & Yes \\
    Year FE & No & Yes & Yes & No \\
    Year-Month FE & No & No & No & Yes \\
    PBU FE & No & No & Yes & Yes \\
    Observations & 196,886 & 196,886 & 196,883 & 196,883 \\
    Within R$^2$ (A) & 0.003 & 0.003 & 0.000 & 0.000 \\
    Within R$^2$ (B) & 0.315 & 0.327 & 0.358 & 0.359 \\
    \hline\hline
  \end{tabular}
  \begin{tablenotes}
    \small
    \item \textit{Notes:} Dependent variable: log negotiated price. Sample: winners only. Panel~A reports the total effect; Panel~B controls for log quantity. Winsorized at 1\%/99\%. Standard errors clustered at the PBU level in parentheses. *** \textit{p}$<$0.01, ** \textit{p}$<$0.05, * \textit{p}$<$0.1.
  \end{tablenotes}
  \end{threeparttable}
\end{table}


\paragraph{Total effect (Panel~A).}
Negotiated prices are significantly higher for urgent purchases across all specifications. The total effect ranges from 17.2\% ($e^{0.159}-1$) with item FE only to 5.4--5.8\% with our preferred specifications (columns~3--4). The attenuation pattern mirrors that of reference prices, confirming the importance of buyer-level heterogeneity.

\paragraph{Direct effect (Panel~B).}
Controlling for log quantity reduces the coefficient on urgency by roughly half in the most saturated specifications: from 0.053 to 0.030 (column~3). The direct effect of 3.1\% ($e^{0.030}-1$, $p < 0.10$) represents the price premium attributable to urgency net of the quantity channel. The coefficient on log quantity is large and negative ($-0.392$), confirming the importance of bulk discounts: a one standard deviation reduction in log quantity is associated with substantially higher per-unit prices.

\subsection{Bidder Participation}
\label{sec:firms}

Table~\ref{tab:firms} examines the effect of urgency on the number of bidding firms.

\begin{table}[ht]
  \centering
  \caption{Participant Firms}
  \label{tab:firms}
  \small
  \begin{threeparttable}
  \begin{tabular}{lcccc}
    \hline\hline
    & (1) & (2) & (3) & (4) \\
    \hline
    \multicolumn{5}{l}{\textit{Panel~A: Total Effect}} \\[3pt]
    Urgent Purchase & -0.101*** & -0.113*** & -0.056*** & -0.054*** \\
     & (0.023) & (0.022) & (0.014) & (0.015) \\[3pt]
    \hline
    \multicolumn{5}{l}{\textit{Panel~B: Direct Effect}} \\[3pt]
    Urgent Purchase & -0.091*** & -0.102*** & -0.042*** & -0.040*** \\
     & (0.014) & (0.011) & (0.012) & (0.013) \\[3pt]
    Log Quantity & 0.063*** & 0.064*** & 0.065*** & 0.065*** \\
     & (0.004) & (0.004) & (0.003) & (0.003) \\[3pt]
    \hline
    Item FE & Yes & Yes & Yes & Yes \\
    Year FE & No & Yes & Yes & No \\
    Year-Month FE & No & No & No & Yes \\
    PBU FE & No & No & Yes & Yes \\
    Obs.\ (Panel~A) & 81,805 & 81,805 & 81,801 & 81,801 \\
    Obs.\ (Panel~B) & 81,793 & 81,793 & 81,789 & 81,789 \\
    Within R$^2$ (A) & 0.004 & 0.006 & 0.001 & 0.001 \\
    Within R$^2$ (B) & 0.063 & 0.069 & 0.043 & 0.044 \\
    \hline\hline
  \end{tabular}
  \begin{tablenotes}
    \small
    \item \textit{Notes:} Dependent variable: log number of bidding firms. Sample: winners only. Panel~A reports the total effect; Panel~B controls for log quantity. Winsorized at 1\%/99\%. Standard errors clustered at the PBU level in parentheses. *** \textit{p}$<$0.01, ** \textit{p}$<$0.05, * \textit{p}$<$0.1.
  \end{tablenotes}
  \end{threeparttable}
\end{table}


\paragraph{Total effect (Panel~A).}
Urgent purchases attract significantly fewer bidding firms across all specifications. In our preferred specification (column~3), urgent tenders have 5.5\% ($e^{-0.056}-1$) fewer participating firms. The total effect reaches 10.7\% with item FE only (column~1).

\paragraph{Direct effect (Panel~B).}
Controlling for quantity only modestly attenuates the coefficient: from $-0.056$ to $-0.042$ in column~(3), corresponding to a direct effect of 4.1\%. The positive coefficient on log quantity (0.065) confirms that larger tenders attract more participants, consistent with standard models of entry into auctions.

The reduction in bidder participation is a key mechanism linking urgency to higher prices. With fewer competing firms, the government's bargaining position weakens. This is consistent with the theoretical prediction that the number of bidders is a primary determinant of auction outcomes \citep{bulow1996auctions}.

\subsection{Tender Success}
\label{sec:success}

Table~\ref{tab:success} reports linear probability model estimates for tender success.

\begin{table}[ht]
  \centering
  \caption{Tender Success (LPM)}
  \label{tab:success}
  \small
  \begin{threeparttable}
  \begin{tabular}{lcccc}
    \hline\hline
    & (1) & (2) & (3) & (4) \\
    \hline
    \multicolumn{5}{l}{\textit{Panel~A: Total Effect}} \\[3pt]
    Urgent Purchase & 0.023** & 0.023** & 0.021*** & 0.021*** \\
     & (0.009) & (0.009) & (0.006) & (0.006) \\[3pt]
    \hline
    \multicolumn{5}{l}{\textit{Panel~B: Direct Effect}} \\[3pt]
    Urgent Purchase & 0.023** & 0.023** & 0.021*** & 0.022*** \\
     & (0.009) & (0.010) & (0.006) & (0.006) \\[3pt]
    Log Quantity & 0.003 & 0.003 & 0.012*** & 0.012*** \\
     & (0.004) & (0.004) & (0.002) & (0.002) \\[3pt]
    \hline
    Item FE & Yes & Yes & Yes & Yes \\
    Year FE & No & Yes & Yes & No \\
    Year-Month FE & No & No & No & Yes \\
    PBU FE & No & No & Yes & Yes \\
    Obs.\ (Panel~A) & 226,305 & 226,305 & 226,301 & 226,301 \\
    Obs.\ (Panel~B) & 226,282 & 226,282 & 226,278 & 226,278 \\
    Within R$^2$ (A) & 0.001 & 0.001 & 0.000 & 0.000 \\
    Within R$^2$ (B) & 0.001 & 0.001 & 0.004 & 0.004 \\
    \hline\hline
  \end{tabular}
  \begin{tablenotes}
    \small
    \item \textit{Notes:} Dependent variable: successful tender (LPM). Sample: all observations. Panel~A reports the total effect; Panel~B controls for log quantity. Winsorized at 1\%/99\%. Standard errors clustered at the PBU level in parentheses. *** \textit{p}$<$0.01, ** \textit{p}$<$0.05, * \textit{p}$<$0.1.
  \end{tablenotes}
  \end{threeparttable}
\end{table}


The coefficient on urgency is positive and statistically significant: urgent purchases are 2.1 percentage points more likely to succeed than ordinary ones (column~3). This result, while perhaps counterintuitive, is consistent with the institutional environment. Under pressure to comply with court orders---and facing severe sanctions for failure---procurement officials may accept higher prices or less favorable terms to ensure that tenders succeed rather than fail. In other words, the observed higher success rate for urgent purchases likely reflects the government's willingness to ``pay the price'' for compliance, consistent with the ``under the gun'' mechanism.

The positive success coefficient also implies that the higher negotiated prices documented in Table~\ref{tab:neg_prices} are not driven by sample selection (if anything, the selection effect works in the opposite direction: more urgent tenders succeed despite less favorable conditions).

\subsection{The ``Under the Gun'' Effect}
\label{sec:utg}

Table~\ref{tab:underthegun} presents estimates comparing litigated and administrative urgent purchases---the core test for the sanction channel.

\begin{table}[ht]
  \centering
  \caption{Under the Gun: Administrative vs Litigated}
  \label{tab:underthegun}
  \small
  \begin{threeparttable}
  \begin{tabular}{lcccc}
    \hline\hline
    & (1) & (2) & (3) & (4) \\
    \hline
    \multicolumn{5}{l}{\textit{Panel~A: Total Effect}} \\[3pt]
    Administrative (vs Litigated) & -0.209*** & -0.205** & -0.262** & -0.265*** \\
     & (0.074) & (0.074) & (0.096) & (0.094) \\[3pt]
    \hline
    \multicolumn{5}{l}{\textit{Panel~B: Direct Effect}} \\[3pt]
    Administrative (vs Litigated) & -0.006 & 0.016 & 0.117 & 0.115 \\
     & (0.056) & (0.058) & (0.122) & (0.115) \\[3pt]
    Log Quantity & -0.284*** & -0.292*** & -0.341*** & -0.344*** \\
     & (0.036) & (0.038) & (0.042) & (0.040) \\[3pt]
    \hline
    Item FE & Yes & Yes & Yes & Yes \\
    Year FE & No & Yes & Yes & No \\
    Year-Month FE & No & No & No & Yes \\
    PBU FE & No & No & Yes & Yes \\
    Observations & 56,803 & 56,803 & 56,803 & 56,802 \\
    Within R$^2$ (A) & 0.011 & 0.010 & 0.007 & 0.007 \\
    Within R$^2$ (B) & 0.287 & 0.292 & 0.307 & 0.310 \\
    \hline\hline
  \end{tabular}
  \begin{tablenotes}
    \small
    \item \textit{Notes:} Dependent variable: log negotiated price. Sample: urgent purchases only (administrative and litigated), winners, items with both types present. Panel~A reports the total effect; Panel~B controls for log quantity. Winsorized at 1\%/99\%. Standard errors clustered at the PBU level in parentheses. *** \textit{p}$<$0.01, ** \textit{p}$<$0.05, * \textit{p}$<$0.1.
  \end{tablenotes}
  \end{threeparttable}
\end{table}


\paragraph{Total effect (Panel~A).}
Administrative purchases have significantly lower negotiated prices than litigated ones across all specifications. In our preferred specification (column~3), the coefficient is $-0.262$ ($p < 0.05$), implying that litigated purchases are 30.0\% ($e^{0.262}-1$) more expensive than comparable administrative purchases. The magnitude is even larger with year-month FE (column~4): 30.3\%.

\paragraph{Direct effect (Panel~B).}
Controlling for log quantity, the coefficient becomes smaller and statistically insignificant (0.117 in column~3). This suggests that a substantial portion of the ``under the gun'' price premium operates through the quantity channel: litigated purchases involve smaller quantities, which translate into higher per-unit prices through reduced bulk discounts. However, the total effect remains economically and statistically significant, reflecting the combined impact of the sanction regime on all procurement dimensions.

These results demonstrate that the threat of judicial sanctions is a quantitatively important driver of procurement costs, beyond the general disruption caused by urgency. The 23--30\% price premium attributable to the sanction channel represents a substantial and previously unquantified cost of judicial enforcement in health procurement. Descriptive geographic evidence on how administrative demand is distributed across municipalities---presented in Appendix~\ref{sec:appendix_admin}---corroborates the institutional parallel between these two procurement channels.

\subsection{Heterogeneity}
\label{sec:heterogeneity}

We examine heterogeneity in the urgency premium along four dimensions. Tables~\ref{tab:het_sus}--\ref{tab:het_pbu} report results from split-sample regressions and interaction models using our preferred specification (item + year + PBU FE).

\begin{table}[ht]
  \centering
  \caption{Heterogeneity: SUS Component}
  \label{tab:het_sus}
  \small
  \begin{threeparttable}
  \begin{tabular}{lccc}
    \hline\hline
    & \multicolumn{2}{c}{Split Sample} & Interaction \\
    \cmidrule(lr){2-3} \cmidrule(lr){4-4}
    & Specialized & Basic~SUS & Full Sample \\
    & (1) & (2) & (3) \\
    \hline
    Urgent Purchase & 0.005 & 0.030* & -0.073 \\
     & (0.046) & (0.016) & (0.097) \\[3pt]
    Urgent × Basic SUS &  &  & 0.159 \\
     &  &  & (0.115) \\[3pt]
    \hline
    FE: Item + Year + PBU & Yes & Yes & Yes \\
    Observations & 45,048 & 151,835 & 196,883 \\
    Within R$^2$ & 0.000 & 0.000 & 0.001 \\
    \hline\hline
  \end{tabular}
  \begin{tablenotes}
    \small
    \item \textit{Notes:} Preferred specification: Item + Year + PBU FE; PBU-clustered SEs in parentheses. Columns (1)--(2) are split-sample estimates. Column (3) reports the interaction. *** \textit{p}$<$0.01, ** \textit{p}$<$0.05, * \textit{p}$<$0.1.
  \end{tablenotes}
  \end{threeparttable}
\end{table}

\begin{table}[ht]
  \centering
  \caption{Heterogeneity: Time Period}
  \label{tab:het_period}
  \small
  \begin{threeparttable}
  \begin{tabular}{lccc}
    \hline\hline
    & \multicolumn{2}{c}{Split Sample} & Interaction \\
    \cmidrule(lr){2-3} \cmidrule(lr){4-4}
    & Early & Late~(2014+) & Full Sample \\
    & (1) & (2) & (3) \\
    \hline
    Urgent Purchase & 0.074*** & 0.045** & 0.259*** \\
     & (0.025) & (0.018) & (0.035) \\[3pt]
    Urgent × Late Period &  &  & -0.339*** \\
     &  &  & (0.043) \\[3pt]
    \hline
    FE: Item + Year + PBU & Yes & Yes & Yes \\
    Observations & 82,658 & 113,936 & 196,883 \\
    Within R$^2$ & 0.000 & 0.000 & 0.005 \\
    \hline\hline
  \end{tabular}
  \begin{tablenotes}
    \small
    \item \textit{Notes:} Preferred specification: Item + Year + PBU FE; PBU-clustered SEs in parentheses. Columns (1)--(2) are split-sample estimates. Column (3) reports the interaction. *** \textit{p}$<$0.01, ** \textit{p}$<$0.05, * \textit{p}$<$0.1.
  \end{tablenotes}
  \end{threeparttable}
\end{table}

\begin{table}[ht]
  \centering
  \caption{Heterogeneity: Market Competition}
  \label{tab:het_competition}
  \small
  \begin{threeparttable}
  \begin{tabular}{lccc}
    \hline\hline
    & \multicolumn{2}{c}{Split Sample} & Interaction \\
    \cmidrule(lr){2-3} \cmidrule(lr){4-4}
    & Low~Comp. & High~Comp. & Full Sample \\
    & (1) & (2) & (3) \\
    \hline
    Urgent Purchase & 0.046*** & 0.143*** & 0.026 \\
     & (0.016) & (0.034) & (0.018) \\[3pt]
    Urgent × High Competition &  &  & 0.226*** \\
     &  &  & (0.044) \\[3pt]
    \hline
    FE: Item + Year + PBU & Yes & Yes & Yes \\
    Observations & 125,950 & 61,313 & 187,264 \\
    Within R$^2$ & 0.000 & 0.001 & 0.001 \\
    \hline\hline
  \end{tabular}
  \begin{tablenotes}
    \small
    \item \textit{Notes:} Preferred specification: Item + Year + PBU FE; PBU-clustered SEs in parentheses. Columns (1)--(2) are split-sample estimates. Column (3) reports the interaction. *** \textit{p}$<$0.01, ** \textit{p}$<$0.05, * \textit{p}$<$0.1.
  \end{tablenotes}
  \end{threeparttable}
\end{table}

\begin{table}[ht]
  \centering
  \caption{Heterogeneity: PBU Size}
  \label{tab:het_pbu}
  \small
  \begin{threeparttable}
  \begin{tabular}{lccc}
    \hline\hline
    & \multicolumn{2}{c}{Split Sample} & Interaction \\
    \cmidrule(lr){2-3} \cmidrule(lr){4-4}
    & Small~PBU & Large~PBU & Full Sample \\
    & (1) & (2) & (3) \\
    \hline
    Urgent Purchase & 0.023 & 0.062*** & -0.000 \\
     & (0.026) & (0.018) & (0.049) \\[3pt]
    Urgent × Large PBU &  &  & 0.070 \\
     &  &  & (0.054) \\[3pt]
    \hline
    FE: Item + Year + PBU & Yes & Yes & Yes \\
    Observations & 31,902 & 164,511 & 196,883 \\
    Within R$^2$ & 0.000 & 0.000 & 0.000 \\
    \hline\hline
  \end{tabular}
  \begin{tablenotes}
    \small
    \item \textit{Notes:} Preferred specification: Item + Year + PBU FE; PBU-clustered SEs in parentheses. Columns (1)--(2) are split-sample estimates. Column (3) reports the interaction. *** \textit{p}$<$0.01, ** \textit{p}$<$0.05, * \textit{p}$<$0.1.
  \end{tablenotes}
  \end{threeparttable}
\end{table}


\paragraph{SUS component (Table~\ref{tab:het_sus}).}
The urgency premium on negotiated prices is concentrated among items in the basic SUS component (primarily common medications): the coefficient is 0.030 ($p < 0.10$) for basic items and an insignificant 0.005 for specialized items. This pattern suggests that the procurement disruption from urgency is more consequential for items where established supply chains and competitive markets normally deliver better prices under ordinary conditions.

\paragraph{Time period (Table~\ref{tab:het_period}).}
The urgency premium is larger in the later period (2014--2019) than in the earlier period (2009--2013), though neither split-sample coefficient is individually significant at the 5\% level. The interaction model reveals a modest differential, consistent with the growing volume and complexity of health litigation over time.

\paragraph{Market competition (Table~\ref{tab:het_competition}).}
Markets with lower competition (below-median number of bidders at the item level) show a somewhat larger urgency premium (0.034 vs.\ 0.029), though the difference is not statistically significant. The pattern is consistent with urgency being most costly when competitive pressure is already weak.

\paragraph{PBU size (Table~\ref{tab:het_pbu}).}
The urgency premium is slightly larger for purchases made by smaller PBUs (below-median transaction volume), consistent with smaller units having less experience and fewer resources to manage urgent procurement effectively.
